\documentclass[11pt,a4paper]{report}

% --- Pacchetti Fondamentali ---
\usepackage[utf8]{inputenc}
\usepackage[T1]{fontenc}
\usepackage{lmodern}
\usepackage[italian]{babel}
\usepackage{etoolbox}
\usepackage[margin=1in]{geometry}
\usepackage{hyperref}
\usepackage{xcolor}
\usepackage{listings}
\usepackage{enumitem}
\usepackage{titlesec}
\usepackage{fancyhdr}
\usepackage{cleveref}
\usepackage{graphicx}
\usepackage{longtable}
\usepackage{booktabs}
\usepackage{setspace}

\crefformat{section}{\S#2#1#3}
\crefformat{subsection}{\S#2#1#3}
\crefformat{subsubsection}{\S#2#1#3}
\crefformat{figure}{#2Figura~#1#3}
\crefformat{footnote}{#2\footnotemark[#1]#3}
\crefformat{table}{#2Tabella~#1#3}

% --- Configurazione Header e Footer ---
\pagestyle{fancy}
\fancyhf{}
\fancyhead[L]{\small Programmazione e Tecnologie Web}
\fancyhead[R]{\small Giuseppe Capasso}
\fancyfoot[C]{\thepage}
\renewcommand{\headrulewidth}{0.4pt}
\usepackage{tikz}
\usetikzlibrary{arrows.meta, positioning, shapes.geometric, decorations.pathreplacing, calc, shadows, fit}

% --- Configurazione Colori e Link ---
\definecolor{primary}{RGB}{38, 66, 139}
\definecolor{codebg}{RGB}{245, 245, 245}
\definecolor{keywordcolor}{RGB}{0, 0, 255}
\definecolor{stringcolor}{RGB}{163, 21, 21}
\definecolor{commentcolor}{RGB}{0, 128, 0}

\hypersetup{
    colorlinks=true,
    linkcolor=primary,
    urlcolor=primary,
    citecolor=primary,
}

% --- Configurazione Formattazione Codice ---
\lstdefinelanguage{json}{
    basicstyle=\ttfamily\small,
    numbers=left,
    numberstyle=\scriptsize,
    stepnumber=1,
    numbersep=8pt,
    showstringspaces=false,
    breaklines=true,
    frame=lines,
    backgroundcolor=\color{codebg},
    string=[s]{"}{"},
    comment=[l]{//},
    morecomment=[s]{/*}{*/},
    stringstyle=\color{stringcolor},
    keywordstyle=\color{keywordcolor},
}

\lstdefinelanguage{javascript}{
    keywords={async, await, break, case, catch, class, const, continue, debugger, default, delete, do, else, export, extends, false, finally, for, function, if, import, in, instanceof, new, null, return, super, switch, this, throw, true, try, typeof, var, void, while, with, yield, let},
    morekeywords={console, log, Error, fetch, JSON, parse, stringify, useState, useEffect, useCallback, useMemo},
    sensitive=true,
    morecomment=[l]{//},
    morecomment=[s]{/*}{*/},
    morestring=[b]',
    morestring=[b]",
    morestring=[b]`,
    keywordstyle=\color{primary}\bfseries,
    commentstyle=\color{commentcolor}\itshape,
    stringstyle=\color{stringcolor},
    basicstyle=\ttfamily\small
}

\lstset{
    backgroundcolor=\color{codebg},
    basicstyle=\ttfamily\small,
    breaklines=true,
    frame=single,
    rulecolor=\color{lightgray},
    captionpos=b,
    xleftmargin=1em,
    xrightmargin=1em,
    keywordstyle=\color{keywordcolor},
    commentstyle=\color{commentcolor},
    stringstyle=\color{stringcolor}
}

% --- Formattazione Titoli ---
\titleformat{\chapter}[display]
  {\normalfont\bfseries\color{primary}}{\huge Capitolo \thechapter}{1ex}{\Huge}
\titlespacing*{\chapter}{0pt}{-20pt}{40pt}

\onehalfspacing

\begin{document}

% --- Frontespizio ---
\begin{titlepage}
    \centering
    \vspace*{4cm}
    {\Huge \textbf{\color{primary} Programmazione e Tecnologie Web}} \\[1.5cm]
    {\Large \textbf{Dispense del Corso - Lezione 5}} \\[0.5cm]
    {\large Docente: Giuseppe Capasso} \\[1.5cm]
    {\normalsize \textbf{Cloud Native Development \& Cybersecurity Specialist}} \\
    {\normalsize Istituto Tecnico Superiore (ITS)} \\
    \vfill
    {\normalsize Anno Accademico 2025/2026}
\end{titlepage}

\newpage
\tableofcontents
\newpage

% ==========================================
% CAPITOLO 1
% ==========================================
\chapter{CSR e la Necessità di un Framework}

\section{Oltre la Manipolazione del DOM}
Nelle lezioni precedenti abbiamo visto come JavaScript possa essere utilizzato per modificare dinamicamente il contenuto di una pagina. Tuttavia, in applicazioni complesse, la gestione manuale del DOM presenta limiti strutturali.

\begin{itemize}
    \item \textbf{Imperativo vs Dichiarativo}: Nello stile imperativo (Vanilla JS), istruiamo il browser passo dopo passo su cosa fare. Nello stile dichiarativo (React), descriviamo \textit{cosa} vogliamo vedere in base allo stato attuale.
    \item \textbf{Sincronizzazione dello Stato}: Gestire manualmente la coerenza tra i dati (es. un array di task) e gli elementi HTML (es. tag \texttt{<li>}) porta inevitabilmente a bug di sincronizzazione.
\end{itemize}

\section{Il Concetto di Virtual DOM}
React non aggiorna l'intera pagina ad ogni cambiamento. Utilizza una tecnica chiamata \textbf{Reconciliation} basata sul \textbf{Virtual DOM}:
\begin{enumerate}
    \item Una copia leggera del DOM reale viene mantenuta in memoria.
    \item Quando i dati cambiano, viene generato un nuovo Virtual DOM.
    \item React calcola la differenza (\textit{Diffing}) tra la vecchia e la nuova versione.
    \item Solo i nodi effettivamente cambiati vengono aggiornati nel browser reale.
\end{enumerate}

\section{Client-Side Rendering (CSR)}
In un'applicazione CSR, il server non invia HTML pronto per la visualizzazione. Invia un file HTML minimo (un "guscio") e un bundle JavaScript. \`E il browser a farsi carico della logica di visualizzazione.

\begin{figure}[ht]
\centering
\begin{tikzpicture}[
    node distance=1.5cm,
    box/.style={rectangle, draw, rounded corners, minimum width=3cm, minimum height=1.2cm, align=center, font=\small},
    arrow/.style={-latex, thick, primary}
]
    \node (browser) [box, fill=blue!5] {Browser\\(Orchestratore)};
    \node (static) [box, fill=gray!10, right=3.5cm of browser] {Static Host / CDN\\(Next.js Bundle)};
    \node (api) [box, fill=primary!20, below=2cm of browser] {Backend API\\(Lezione 4)};

    \draw [arrow] ([yshift=0.2cm]browser.east) -- node[above, scale=0.7] {1. Richiesta Bundle} ([yshift=0.2cm]static.west);
    \draw [arrow] ([yshift=-0.2cm]static.west) -- node[below, scale=0.7] {2. Esecuzione JS} ([yshift=-0.2cm]browser.east);
    
    \draw [arrow] ([xshift=-0.2cm]browser.south) -- node[left, scale=0.7] {3. Fetch JSON} ([xshift=-0.2cm]api.north);
    \draw [arrow] ([xshift=0.2cm]api.north) -- node[right, scale=0.7] {4. Renderizza Dati} ([xshift=0.2cm]browser.south);

    \node at ($(browser.north)+(0,0.5cm)$) [font=\bfseries\scriptsize, primary] {Hydration \& Reattività};
\end{tikzpicture}
\caption{Ciclo di vita CSR: il browser agisce come hub centrale per risorse e dati.}
\label{fig:csr_flow}
\end{figure}

\newpage

% ==========================================
% CAPITOLO 2 (NUOVO)
% ==========================================
\chapter{Sicurezza e Comunicazione: CORS e Cookie}

\section{La Same-Origin Policy (SOP)}
Il browser implementa una restrizione di sicurezza fondamentale chiamata \textbf{Same-Origin Policy}. Essa impedisce a uno script caricato da un'origine (es. \texttt{http://localhost:3000}) di interagire con risorse di un'altra origine (es. \texttt{http://localhost:5000}).
Due URL hanno la stessa origine se condividono esattamente: \textbf{Protocollo}, \textbf{Dominio} e \textbf{Porta}.

\section{CORS: Cross-Origin Resource Sharing}
Quando sviluppiamo un frontend React (porta 3000) che interroga un backend Flask (porta 5000), violiamo la SOP. Il meccanismo **CORS** permette al server di "autorizzare" esplicitamente il browser a leggere la risposta.

\subsection{Il Preflight Request (OPTIONS)}
Per richieste potenzialmente pericolose (es. con header personalizzati o metodi come DELETE), il browser invia prima una richiesta automatica \textbf{OPTIONS}. Se il server risponde con gli header corretti, la richiesta reale viene inviata.

\begin{itemize}
    \item \texttt{Access-Control-Allow-Origin}: Specifica quali domini possono accedere.
    \item \texttt{Access-Control-Allow-Methods}: Metodi consentiti (GET, POST, ecc).
\end{itemize}

\section{Cookie e Gestione dello Stato}
I cookie sono piccoli frammenti di dati che il server invia al browser e che il browser memorizza, rispedendoli automaticamente al server ad ogni richiesta successiva verso quel dominio.

\subsection{Attributi di Sicurezza}
\begin{itemize}
    \item \textbf{HttpOnly}: Impedisce a JavaScript di leggere il cookie (protezione contro XSS).
    \item \textbf{Secure}: Il cookie viene inviato solo su connessioni HTTPS.
    \item \textbf{SameSite}: Controlla se il cookie deve essere inviato durante richieste cross-site (Lax, Strict, None).
\end{itemize}

\newpage

% ==========================================
% CAPITOLO 3
% ==========================================
\chapter{Approccio Monolitico in React}

\section{Il Componente "Tutto-in-Uno"}
Per comprendere i vantaggi di React, iniziamo costruendo la Todo App in un singolo file. Questo approccio è utile per piccoli prototipi, ma mostra rapidamente i suoi limiti di leggibilità.

\begin{lstlisting}[language=javascript, caption=Todo App Monolitica (app/page.js)]
"use client";
import { useState, useEffect } from "react";

export default function MonolithicTodo() {
  const [tasks, setTasks] = useState([]);
  const [search, setSearch] = useState("");

  // 1. Data Fetching all'avvio
  useEffect(() => {
    fetch("http://localhost:3000/tasks")
      .then(res => res.json())
      .then(data => setTasks(data));
  }, []);

  // 2. Stato derivato per il filtraggio
  const filteredTasks = tasks.filter(t => 
    t.text.toLowerCase().includes(search.toLowerCase())
  );

  return (
    <main className="p-10">
      <input 
        value={search} 
        onChange={(e) => setSearch(e.target.value)}
        placeholder="Filtra task..."
        className="border p-2 w-full mb-4"
      />
      <ul className="space-y-2">
        {filteredTasks.map((t, index) => (
          <li key={index} className="p-2 bg-white shadow rounded flex justify-between">
            {t.text}
            <button onClick={() => setTasks(tasks.filter((_, i) => i !== index))}>
              X
            </button>
          </li>
        ))}
      </ul>
    </main>
  );
}
\end{lstlisting}

\section{Il Ciclo di Vita di un Componente}
Il componente attraversa diverse fasi. Comprendere dove si collocano \texttt{useState} e \texttt{useEffect} è fondamentale.

\begin{figure}[ht]
\centering
\begin{tikzpicture}[
    node distance=1cm,
    phase/.style={rectangle, draw, rounded corners, minimum width=3.5cm, minimum height=1cm, align=center, font=\small},
    arrow/.style={-latex, thick, primary}
]
    \node (mount) [phase, fill=green!10] {\textbf{Mounting}\\(Nascita)};
    \node (render) [phase, below=of mount, fill=blue!5] {Render\\(JSX -> VDOM)};
    \node (effect) [phase, right=2cm of render, fill=yellow!10] {Effect\\(API, Timer, ecc)};
    \node (update) [phase, below=of render, fill=orange!10] {\textbf{Updating}\\(Cambiamento Stato)};
    \node (unmount) [phase, below=of update, fill=red!10] {\textbf{Unmounting}\\(Morte)};

    \draw [arrow] (mount) -- (render);
    \draw [arrow] (render) -- node[above, scale=0.7] {Completato} (effect);
    \draw [arrow] (effect) -- ++(1cm,0) |- node[near start, right, scale=0.7] {Cleanup} (update);
    \draw [arrow] (update) -- (render);
    \draw [arrow] (update) -- (unmount);
\end{tikzpicture}
\caption{Diagramma del ciclo di vita: l'effetto viene eseguito dopo il rendering visuale.}
\label{fig:react_lifecycle}
\end{figure}

\newpage

% ==========================================
% CAPITOLO 4
% ==========================================
\chapter{Rifattorizzazione e Best Practices}

\section{Astrazione in Componenti}
Man mano che l'app cresce, il file monolitico diventa difficile da gestire. Applichiamo il principio di singola responsabilità scomponendo la UI.

\begin{lstlisting}[language=javascript, caption=Componenti separati]
// Sotto-componente atomico
const TaskItem = ({ text, onDelete }) => (
  <li className="flex justify-between p-2 border-b">
    <span>{text}</span>
    <button onClick={onDelete} className="text-red-500">Elimina</button>
  </li>
);

// Componente di visualizzazione lista
const TodoList = ({ data, onItemDelete }) => (
  <ul className="mt-4">
    {data.map((t, i) => (
      <TaskItem key={i} text={t.text} onDelete={() => onItemDelete(i)} />
    ))}
  </ul>
);
\end{lstlisting}

\section{Ottimizzazione con \texttt{useCallback}}
Quando passiamo funzioni ai figli (\textit{Props Drilling}), queste vengono ricreate ad ogni render del genitore. Usiamo \texttt{useCallback} per memorizzarle.

\begin{lstlisting}[language=javascript]
const removeTask = useCallback((index) => {
  setTasks(prev => prev.filter((_, i) => i !== index));
}, []); // La funzione non viene mai ricreata
\end{lstlisting}

\section{API Routes: Verso l'Unified Full-Stack}
Next.js permette di scrivere logica backend in file \texttt{route.js}. Questo elimina i problemi di CORS poich\'e frontend e backend condividono la stessa origine.

\begin{lstlisting}[language=javascript, caption=Esempio API Route (app/api/stats/route.js)]
export async function GET() {
  return Response.json({ 
    serverTime: new Date().toISOString(),
    status: "online" 
  });
}
\end{lstlisting}

\section{Conclusioni}
Siamo passati da un componente disordinato a una struttura professionale, scalabile e performante. React ci insegna che il codice pulito nasce dalla corretta gestione dei dati e dalla scomposizione logica dell'interfaccia.

\end{document}
